% Source PDF: source/普通高考/2019/2019全国3理(云南、广西、贵州、西藏、四川).pdf, page 1.
% PDF embedded bitmap attempt: pdfimages -list found no embedded images; the source histogram is vector page content.
% Grid evidence: tmp/repaint_2019_national_paper_3_science/page1-grid100.png (major 100px, fine 50px).
% No-grid crop evidence: tmp/repaint_2019_national_paper_3_science/q17_fig2_orig_try.png.
% Elements: frequency/group-width axis with arrow, percentage axis with break, six bars, dashed reference levels, labels and caption.
% Relations: bars occupy compressed intervals 2.5--3.5 through 7.5--8.5 with frequencies 0.05, b, 0.15, a, 0.20, 0.15; the axis break compresses 0--2.5.
% solid segments: axes, arrowheads, tick marks, bar outlines, axis-break zigzag.
% dashed segments: horizontal guides at 0.05, b, 0.15, 0.20, and a ending at the first bar boundary for each level.
% labels: frequency/group-width above the y-axis, percentage to the right of the x-axis, a/b and numeric ticks, Chinese caption below.
\begin{tikzpicture}[baseline]
\begin{axis}[scaled ticks=false,exam statistical histogram,
    width=15.806cm,
    height=9.031cm,
    axis lines=none,
    xmin=-0.55,
    xmax=8.15,
    ymin=-0.12,
    ymax=0.36,
    clip=false,
    ticks=none,
]
    \draw[dashed,line width=0.55pt] (axis cs:0,0.05) -- (axis cs:1,0.05);
    \draw[dashed,line width=0.55pt] (axis cs:0,0.10) -- (axis cs:2,0.10);
    \draw[dashed,line width=0.55pt] (axis cs:0,0.15) -- (axis cs:3,0.15);
    \draw[dashed,line width=0.55pt] (axis cs:0,0.20) -- (axis cs:5,0.20);
    \draw[dashed,line width=0.55pt] (axis cs:0,0.30) -- (axis cs:4,0.30);

    \addplot[
        ybar interval,
        draw=black,
        fill=white,
        line width=0.55pt,
    ] coordinates {
        (1,0.05)
        (2,0.10)
        (3,0.15)
        (4,0.30)
        (5,0.20)
        (6,0.15)
        (7,0.00)
    };

    \draw[->,line width=0.65pt] (axis cs:0,0) -- (axis cs:8.0,0);
    \draw[->,line width=0.65pt] (axis cs:0,0) -- (axis cs:0,0.345);

    % The source axis has a break immediately to the right of the origin.
    \draw[white,line width=2.5pt] (axis cs:0.15,0) -- (axis cs:0.70,0);
    \examstatxbreak[line width=0.55pt]{0.405}{0};

    \draw[line width=0.4pt] (axis cs:0,0) -- (axis cs:0,-0.012);
    \draw[line width=0.4pt] (axis cs:1,0) -- (axis cs:1,-0.012);
    \draw[line width=0.4pt] (axis cs:2,0) -- (axis cs:2,-0.012);
    \draw[line width=0.4pt] (axis cs:3,0) -- (axis cs:3,-0.012);
    \draw[line width=0.4pt] (axis cs:4,0) -- (axis cs:4,-0.012);
    \draw[line width=0.4pt] (axis cs:5,0) -- (axis cs:5,-0.012);
    \draw[line width=0.4pt] (axis cs:6,0) -- (axis cs:6,-0.012);
    \draw[line width=0.4pt] (axis cs:7,0) -- (axis cs:7,-0.012);
    \node[anchor=north,inner sep=0pt,font=\normalsize,yshift=-2pt] at (axis cs:0,-0.012000) {0};
    \node[anchor=north,inner sep=0pt,font=\normalsize,yshift=-2pt] at (axis cs:1,-0.012000) {2.5};
    \node[anchor=north,inner sep=0pt,font=\normalsize,yshift=-2pt] at (axis cs:2,-0.012000) {3.5};
    \node[anchor=north,inner sep=0pt,font=\normalsize,yshift=-2pt] at (axis cs:3,-0.012000) {4.5};
    \node[anchor=north,inner sep=0pt,font=\normalsize,yshift=-2pt] at (axis cs:4,-0.012000) {5.5};
    \node[anchor=north,inner sep=0pt,font=\normalsize,yshift=-2pt] at (axis cs:5,-0.012000) {6.5};
    \node[anchor=north,inner sep=0pt,font=\normalsize,yshift=-2pt] at (axis cs:6,-0.012000) {7.5};
    \node[anchor=north,inner sep=0pt,font=\normalsize,yshift=-2pt] at (axis cs:7,-0.012000) {8.5};
    \draw[line width=0.4pt] (axis cs:0,0.05) -- (axis cs:-0.10,0.05);
    \draw[line width=0.4pt] (axis cs:0,0.15) -- (axis cs:-0.10,0.15);
    \draw[line width=0.4pt] (axis cs:0,0.20) -- (axis cs:-0.10,0.20);
    \node[anchor=east,inner sep=0pt,font=\normalsize,xshift=-4pt] at (axis cs:-0.100000,0.05) {0.05};
    \node[anchor=east,inner sep=0pt,font=\normalsize,xshift=-4pt] at (axis cs:-0.100000,0.15) {0.15};
    \node[anchor=east,inner sep=0pt,font=\normalsize,xshift=-4pt] at (axis cs:-0.100000,0.20) {0.20};
    \draw[line width=0.4pt] (axis cs:0,0.10) -- (axis cs:-0.10,0.10);
    \node[anchor=east,inner sep=0pt,font=\normalsize,xshift=-4pt] at (axis cs:-0.05,0.10) {$b$};
    \draw[line width=0.4pt] (axis cs:0,0.30) -- (axis cs:-0.10,0.30);
    \node[anchor=east,inner sep=0pt,font=\normalsize,xshift=-4pt] at (axis cs:-0.05,0.30) {$a$};

    \examstatylabel{\(\dfrac{\text{频率}}{\text{组距}}\)};
    \examstatxlabel{百分比};
    \node[anchor=north,inner sep=0pt,font=\normalsize,yshift=-4pt] at (axis cs:4.0,-0.09) {乙离子残留百分比直方图};
    \path[use as bounding box] (axis cs:-0.80,-0.12) rectangle (axis cs:8.15,0.36);
\end{axis}
\end{tikzpicture}
